\section{Pricing the rental-default insurance}\label{sec:insurance}

Quoted French rental insurance (\textit{garantie loyers impayés}) is
typically priced at around 2--3\% of the gross annual rent. We benchmark
two consistent ways of valuing it endogenously: the actuarial
standard-deviation principle and an option-theoretic put on rent
shortfall.

\subsection{Actuarial price}

Following the classical \autocite{EmbrechtsMcNeilStraumann2002} treatment,
the insurer absorbs a fraction $\kappa$ of every realised default loss and
charges a premium
\begin{equation}
P_\text{act} = (1 + \theta)\, \E[L_\text{covered}] + \lambda\,
\sigma(L_\text{covered}),
\label{eq:act-premium}
\end{equation}
where $L_\text{covered} = \kappa \cdot L(T) \cdot P_0$ is the covered
lifetime loss per portfolio path. We use the configured loadings
$\theta = 0.10$ (admin) and $\lambda = 0.15$ (risk) and coverage $\kappa =
\insuranceCoverage{}$. The premium can be paid as a lump sum at $t = 0$
or, equivalently in PV terms, as a level annual payment over the horizon
$T$. On the Paris-intermediate baseline the expected covered loss is
$\E[L_\text{covered}] \approx \insuranceExpectedLoss{}$, which translates
into an annual-equivalent premium of \insuranceAnnualPremium{}.

\subsection{Option-theoretic price}

Alternatively, insurance is a put on the periodic net rent. Let
$R^{\text{prom}}_t$ denote the promised rent (full-tenancy rate net of a
one-month deductible) and $R_t$ the realised flow. The insurer pays
$\max(R^{\text{prom}}_t - R_t, 0)$ each period; under the simulated short
rate $r_t$ the lump-sum premium is
\begin{equation}
P_\text{opt} = \E\left[\sum_{k=1}^{N} D(t_k)\,
\max\bigl(R^{\text{prom}}_{t_k} - R_{t_k}, 0\bigr) \cdot \kappa
\right],
\qquad D(t_k) = \exp\!\Big(-\textstyle\int_0^{t_k} r_s\,\mathrm{d}s\Big).
\label{eq:opt-premium}
\end{equation}

\subsection{Tranche-level allocation}\label{ssec:tranche-allocation}

A first-pass implementation would subtract a flat annual premium from
the aggregate net rent stream. That choice mis-allocates the burden
across tranches: the senior, which absorbs almost no default loss in
practice, would pay the same premium drag as the equity, which absorbs
most of it. We instead allocate the premium proportional to each
tranche's \emph{expected covered loss} under the no-insurance baseline:
\begin{equation}
w_\tau = \frac{\E[L_\tau(L(T))] \cdot \kappa}{\sum_{\tau'} \E[L_{\tau'}(L(T))] \cdot \kappa},
\qquad \text{premium drag on } \tau = w_\tau \cdot P_{\text{act}} \cdot \Delta t.
\label{eq:premium-allocation}
\end{equation}
With the calibrated parameters and the stop-loss attach/detach scheme
$(0, 0.25), (0.25, 0.60), (0.60, 1)$ the equity tranche carries roughly
$w_\text{equity} \approx 84\,\%$ of the premium, the mezzanine
$\approx 16\,\%$, and the senior less than $1\,\%$ --- the
seniority-aware allocation a fair pricing requires.

\subsection{Side-by-side comparison on the calibrated baseline}\label{ssec:insurance-comparison}

The two pricing approaches produce internally consistent but
quantitatively different premiums on our Paris baseline:
the standard-deviation principle bundles an explicit risk loading on
top of the expected covered loss, while the option-theoretic price
discounts only the realised shortfall under the simulated short rate.
On the same MC paths the option-theoretic premium therefore sits
\emph{below} the actuarial one. \Cref{tab:insurance-comparison} and
\cref{fig:insurance-comparison} report the per-instrument fair-to-par
under the three regimes.


\begin{table}[ht]
\centering
\caption{Insurance regime comparison: fair price / par per instrument with no insurance, actuarial pricing and option-theoretic pricing (Gaussian one-factor copula baseline).}
\label{tab:insurance-comparison}
\small
\begin{tabular}{l r r r}
\toprule
Instrument & No insurance & Actuarial & Option-theor. \\
\midrule
Model A & 1.107 & 0.909 & 0.969 \\
Model B & 0.908 & 0.905 & 0.964 \\
Equity & 1.030 & 1.044 & 1.243 \\
Mezzanine & 0.786 & 0.766 & 0.794 \\
Senior & 0.939 & 0.939 & 0.939 \\
\bottomrule
\end{tabular}
\end{table}


\begin{figure}[ht]
  \centering
  \includegraphics[width=\linewidth]{fig_insurance_comparison.pdf}
  \caption{Grouped bars of fair price / par across no insurance,
  actuarial insurance, and option-theoretic insurance, under the
  Gaussian one-factor copula baseline.}
  \label{fig:insurance-comparison}
\end{figure}

Equity benefits the most from coverage under both pricing regimes —
the structural cushion at $[0, 25\,\%]$ absorbs the bulk of the
realised loss, so insuring 90\,\% of it sharply lifts its fair-to-par.
The option-theoretic premium delivers a larger lift than the actuarial
one because the actuarial loading $\lambda \cdot \sigma(L_{\mathrm{covered}})$
is what investors pay for the systemic-risk hedge.

\subsection{Discussion}

The two prices coincide only for specific loadings; they differ here
because the standard-deviation principle charges for total-loss
variance whereas the option-theoretic price charges only for the
realised shortfall. The gap quantifies the implicit risk premium the
actuarial principle attaches to the systemic factor. Both prices are
written to \texttt{artifacts/results\_meta.json}; the comparison table
sits in \cref{sec:results}.
